\documentclass[10pt,a4paper]{amsart}
\usepackage[margin=19mm]{geometry}
\usepackage{amsmath,amssymb,mathtools}
\usepackage[scr=rsfs,cal=euler]{mathalfa}
\usepackage{xfrac}
\usepackage[unicode,hidelinks,bookmarks=false]{hyperref}
\newcommand{\ie}{i.e.\ }
\newcommand{\cf}{cf.\ }
\newcommand{\resp}{respectively\ }
\newcommand{\Subsection}{\subsection}
\providecommand{\nobreakdash}{\nobreak\hbox{-}}
\setlength{\emergencystretch}{2em}
\multlinegap0pt
\usepackage[shortlabels]{enumitem}
\usepackage{mathtools}
\usepackage{tikz}
\newtheorem{lemma}{Lemma}
\title{A counterexample to the bounded mass property}
\author{Mingchen Xia, Kewei Zhang}
\date{Source: arXiv:2608.21053v1 (CC BY 4.0)}
\begin{document}
\maketitle

\noindent\emph{Source and licence.} A counterexample to the bounded mass property. Version arXiv:2608.21053v1 (CC BY 4.0), distributed by arXiv under the Creative Commons Attribution 4.0 International licence, http://creativecommons.org/licenses/by/4.0/, as recorded in the arXiv licence metadata for this exact version and shown on the abstract page https://arxiv.org/abs/2608.21053v1. Full licence notice: \texttt{licenses/arxiv-2608.21053-CC-BY-4.0.txt}.\par\smallskip

\noindent\textit{Editorial scope.} Counted unit: the article's lemma lem:heat-fiber (source0.tex lines 562--605). The article's complete original proof of it is delivered with it (source0.tex lines 607--701, one \texttt{proof} environment of 95 lines). No mathematics is written, repaired or re-derived: the statement and the proof are the article's own lines, in the article's own notation, and no retained line is substituted or re-flowed.  Retained uncounted as the source-local context the proof needs: lines 488--494; lines 525--534; lines 546--550.  Nothing was omitted from any retained span, so the package carries no omission record.  No retained line contains a citation, so the wrapper carries no bibliography and no bibliography entry.  No retained line contains a figure environment, an image-inclusion command or drawing code, so no image is delivered and the package declares no asset dependency.  Editorial formatting only, in addition: an AMS \texttt{amsart} preamble that restates the article's own declarations and macros used by the retained lines, namely the environments \texttt{lemma} on separate counters, together with the standard packages those lines need.  The retained environments print in this document as Lemma 1 in order of appearance, whereas the article prints the same items with the numbering of its own full document.  The complete native source of the article as distributed in the arXiv e-print for this exact version is bundled unchanged as \texttt{sources/208221/source0.tex}.
	\begin{equation}
		k_t(w)\coloneqq \frac{2}{t}\sum_{\lambda\in\Lambda}
		\exp\left(-\frac{|w-\lambda|^2}{t}\right),
		\qquad
		\Lambda:=\mathbb Z+2\pi\mathrm{i}\mathbb Z,
		\label{eq:heat-kernel}
	\end{equation}

	\begin{equation}
		\partial_w\chi_{p,q}
		=\left(\pi\mathrm{i}p+\frac q2\right)\chi_{p,q},
		\qquad
		\partial_{\bar w}\chi_{p,q}
		=\left(\pi\mathrm{i}p-\frac q2\right)\chi_{p,q},
		\qquad
		(\chi_{p,q})_{w\bar w}=-\pi^2|\nu_{p,q}|^2\,\chi_{p,q}.
		\label{eq:character-derivatives}
	\end{equation}

	\begin{equation}
		k_t(w)=\sum_{(p,q)\in\mathbb Z^2}
		\mathrm{e}^{-\pi^2|\nu_{p,q}|^2t}\,\chi_{p,q}(w),
		\label{eq:heat-Fourier}
	\end{equation}

	\begin{lemma}
		\label{lem:heat-fiber}
		For $0<\eta<1/10$, define
		$F_\eta\colon(0,\infty)\times E\to\mathbb R$ by
		\begin{equation}
			F_\eta(t,w)\coloneqq -\eta\int_t^\infty\bigl(k_r(w)-1\bigr)\,\mathrm{d}r
			=-\frac{\eta}{\pi^2}
			\sum_{(p,q)\ne(0,0)}
			\frac{\mathrm{e}^{-\pi^2|\nu_{p,q}|^2t}}
			{|\nu_{p,q}|^2}\,\chi_{p,q}(w).
			\label{eq:heat-potential}
		\end{equation}
		Then $F_\eta$ is smooth and real, and the following hold.
		\begin{enumerate}[label={(\arabic*)}]
			\item One has
			\begin{equation}
				(F_\eta)_t=(F_\eta)_{w\bar w}=\eta(k_t-1),
				\qquad
				1+(F_\eta)_{w\bar w}=1-\eta+\eta k_t>1-\eta.
				\label{eq:heat-identity}
			\end{equation}
			\item One has
			\begin{equation}
				(\log k_t)_{w\bar w}\ge-\frac1t
				\qquad\text{on }E,\quad t>0.
				\label{eq:log-heat}
			\end{equation}
			\item For $0<t\le1$,
			\begin{equation}
				\int_E|(F_\eta)_w(t,w)|^2\,\mathrm{d}A_E
				=2\eta^2\log\frac1t+O(\eta^2).
				\label{eq:heat-energy}
			\end{equation}
			Here the implicit constant is uniform for $0<t\le1$ and
			$0<\eta<1/10$.
			\item $F_\eta$ is exponentially small in the sense that, for all fixed integers $a,m\ge0$ and
			all $t\ge1$:
			\begin{equation}
				\|\partial_t^a F_\eta(t,\cdot)\|_{C^m(E)}
				\le C_{a,m}\,\eta\,\mathrm{e}^{-t/4}.
				\label{eq:heat-flat-F}
			\end{equation}
		\end{enumerate}
	\end{lemma}

	\begin{proof}
		Since
		$k_r-1\to0$ exponentially fast, the integral in
		\eqref{eq:heat-potential} converges, and termwise integration in $r$
		gives the stated Fourier series.  The series and all its
		$(t,w)$-derivatives converge uniformly on $[t_0,\infty)\times E$ for
		each $t_0>0$, so $F_\eta$ is smooth; it is real because the
		coefficients are real and invariant under
		$(p,q)\mapsto(-p,-q)$.
		
		(1)  Differentiating \eqref{eq:heat-potential} in $t$ gives
		$(F_\eta)_t=\eta(k_t-1)$; differentiating the Fourier series in
		$w,\bar w$ and using \eqref{eq:character-derivatives} gives
		\[
		(F_\eta)_{w\bar w}
		=\eta\sum_{(p,q)\ne(0,0)}\mathrm{e}^{-\pi^2|\nu_{p,q}|^2t}\chi_{p,q}
		=\eta(k_t-1). 
		\]
		The last inequality in \eqref{eq:heat-identity} holds
		because $k_t>0$.
		
		(2)  Fix a lift $w$ and introduce the probability weights
		\[
		P_{\lambda}\coloneqq \frac{\mathrm{e}^{-|w-\lambda|^2/t}}
		{\sum_{\lambda'\in\Lambda}\mathrm{e}^{-|w-\lambda'|^2/t}},
		\qquad\lambda\in\Lambda .
		\]
		From $\partial_{\bar w}|w-\lambda|^2=w-\lambda$ we get
		\begin{equation}
			t\,\partial_{\bar w}\log k_t(w)
			=-\sum_{\lambda\in\Lambda}(w-\lambda)\,P_\lambda ,
			\label{eq:heat-log-gradient-exact}
		\end{equation}
		and differentiating once more, now in $w$, we find
		\[
		(\log k_t)_{w\bar w}
		=-\frac1t+\frac1{t^2}
		\left\{
		\sum_{\lambda}|w-\lambda|^2P_{\lambda}
		-\Bigl|\sum_{\lambda}(w-\lambda)P_{\lambda}\Bigr|^2
		\right\}.
		\]
		The bracket is the variance of the $\mathbb C$-valued random variable
		$w-\lambda$ under the probability distribution $(P_\lambda)$, hence
		nonnegative.  This proves \eqref{eq:log-heat}.
		
		(3) Parseval's identity applied to \eqref{eq:heat-potential} yields
		\begin{equation}
			\mathcal E_\eta(t)\coloneqq \int_E|(F_\eta)_w|^2\,\mathrm{d}A_E
			=\frac{\eta^2}{\pi^2}
			\sum_{(p,q)\ne(0,0)}
			\frac{\mathrm{e}^{-2\pi^2|\nu_{p,q}|^2t}}
			{|\nu_{p,q}|^2}.
			\label{eq:heat-parseval}
		\end{equation}
		Differentiating termwise and comparing with \eqref{eq:heat-Fourier}, we find
		\[
		\mathcal E_\eta'(t)
		=-2\eta^2\sum_{(p,q)\ne(0,0)}\mathrm{e}^{-2\pi^2|\nu_{p,q}|^2t}
		=-2\eta^2\bigl(k_{2t}(0)-1\bigr).
		\]
		On the other hand, evaluating \eqref{eq:heat-kernel} at $w=0$ and time $2r$, and using $|\lambda|\ge1$ together with the elementary inequality 
		\(
		\mathrm{e}^{-|\lambda|^2/(2r)}\le
		\mathrm{e}^{-1/(4r)}\mathrm{e}^{-|\lambda|^2/4}
		\)
		for $0<r\le1$, we get
		\[
		k_{2r}(0)
		=\frac1r\Bigl(1+
		\sum_{\lambda\in\Lambda\setminus\{0\}}
		\mathrm{e}^{-|\lambda|^2/(2r)}\Bigr)
		=\frac1r+O\left(\frac{\mathrm{e}^{-1/(4r)}}{r}\right),
		\qquad 0<r\le1 .
		\]
		The error $r^{-1}\mathrm{e}^{-1/(4r)}$ is integrable on $(0,1]$, and
		$\mathcal E_\eta(1)\le C\eta^2$ by \eqref{eq:heat-parseval}.
		Integrating $\mathcal E_\eta'$ over $[t,1]$ therefore gives
		\eqref{eq:heat-energy}.
		
		(4)  Fix $a$ and $m$, and abbreviate $\nu\coloneqq\nu_{p,q}$.  Every
		derivative $\partial_t^a\partial^{m'}$ with $m'\le m$, applied to the
		series \eqref{eq:heat-potential}, is a sum over $(p,q)\ne(0,0)$ of terms
		bounded by a polynomial in $|\nu|$ times
		$\eta\mathrm{e}^{-\pi^2|\nu|^2t}$.  Writing
		\[
		\mathrm{e}^{-\pi^2|\nu|^2t}
		=\mathrm{e}^{-t/4}\,\mathrm{e}^{-(\pi^2|\nu|^2-1/4)t}
		\le \mathrm{e}^{-t/4}\,\mathrm{e}^{-(\pi^2|\nu|^2-1/4)}
		\]
		for $t\ge1$, and noting that
		$\sum_{\nu\ne 0}\mathrm{polynomial}(|\nu|)\,
		\mathrm{e}^{-(\pi^2|\nu|^2-1/4)}<\infty$, we obtain
		\eqref{eq:heat-flat-F}.
	\end{proof}

\end{document}
